\documentclass[10pt,twocolumn]{article}

\usepackage[utf8]{inputenc}
\usepackage[T1]{fontenc}
\usepackage{geometry}
\usepackage{amsmath}
\usepackage{amssymb}
\usepackage{booktabs}
\usepackage{balance}
\usepackage{graphicx}
\usepackage{authblk}
\usepackage{abstract}
\usepackage{silence}
\WarningFilter{latex}{Command \showhyphens has changed}
\usepackage{microtype}
\usepackage[hidelinks]{hyperref}
\usepackage[
  backend=biber,
  style=authoryear,
  natbib=true,
  maxbibnames=10,
  minbibnames=10,
  maxcitenames=2,
  giveninits=true,
  uniquename=init,
  doi=true,
  url=false,
  isbn=false,
  eprint=false
]{biblatex}
\addbibresource{refs.bib}
\usepackage{titlesec}
\usepackage{caption}
\usepackage{url}
\usepackage{fancyhdr}
\usepackage{lastpage}
\usepackage{xcolor}
\usepackage{newtxtext,newtxmath}

\hbadness=10000
\vbadness=10000

\geometry{margin=0.75in}
\setlength{\columnsep}{0.25in}
\renewcommand\Authfont{\normalsize\bfseries}
\renewcommand\Affilfont{\small\itshape}
\setlength{\affilsep}{0.4em}
\renewcommand{\abstractname}{}
\setlength{\absleftindent}{0.4in}
\setlength{\absrightindent}{0.4in}
\setlength{\abstitleskip}{0pt}
\titlespacing*{\section}{0pt}{1.1em}{0.4em}
\titlespacing*{\subsection}{0pt}{0.9em}{0.3em}
\titleformat{\section}{\large\bfseries}{\thesection.}{0.4em}{}
\titleformat{\subsection}{\normalsize\bfseries}{\thesubsection.}{0.4em}{}
\captionsetup{font=small,labelfont=bf}
\setcounter{topnumber}{2}
\setcounter{dbltopnumber}{2}
\renewcommand{\topfraction}{0.9}
\renewcommand{\dbltopfraction}{0.9}
\renewcommand{\textfraction}{0.05}
\renewcommand{\floatpagefraction}{0.8}
\renewcommand{\dblfloatpagefraction}{0.8}

\newcommand{\leadauthor}{Swarna}
\newcommand{\shorttitle}{clrcycle for cyclic transcriptomic structure}
\newcommand{\biorxivwordmark}{bio\textcolor{red}{R}$\chi$iv}

\pagestyle{fancy}
\fancyhf{}
\renewcommand{\headrulewidth}{0pt}
\renewcommand{\footrulewidth}{0.4pt}
\fancyfoot[R]{\footnotesize\itshape \leadauthor\ et al. \hspace{6pt}|\hspace{6pt} \biorxivwordmark \hspace{6pt}|\hspace{6pt} \thepage--\pageref{LastPage}}
\fancypagestyle{firststyle}{
  \fancyhf{}
  \renewcommand{\headrulewidth}{0pt}
  \renewcommand{\footrulewidth}{0.4pt}
  \fancyfoot[R]{\footnotesize\itshape \leadauthor\ et al. \hspace{6pt}|\hspace{6pt} \biorxivwordmark \hspace{6pt}|\hspace{6pt} \today \hspace{6pt}|\hspace{6pt} \thepage--\pageref{LastPage}}
}

\title{Circular Projection for Cyclic Compositional Data}
\author[1]{Nikhila Swarna}
\author[1]{Jawhara Emhemed}
\author[1,*]{Lior Pachter}
\affil[1]{Division of Biology and Biological Engineering, California Institute of Technology, Pasadena, CA}
\affil[*]{Address correspondence to: \href{mailto:lpachter@caltech.edu}{lpachter@caltech.edu}}
\date{}

\begin{document}
\thispagestyle{firststyle}

\twocolumn[
\maketitle
\vspace{-1.5em}

\begin{onecolabstract}
Transcriptomes often vary along periodic rather than linear coordinates.
Standard displays can reveal such variation but do not link a sample's position
on a cycle to an ordering of the molecular features that define it. clrcycle
applies a centered log-ratio transformation, learns a cyclic feature ordering by
maximizing first-Fourier-mode variance, and projects samples onto the corresponding
cosine and sine coordinates. Each sample receives a phase and a radius measuring
the strength of its first-harmonic compositional signal. In the Hogenesch mouse
circadian atlas, a supervised reference is contrasted with label-free selection
by two-cycle reproducibility, which recovers phase in ten of twelve tissues and
reveals tissue-specific changes in rhythmic program strength.
\end{onecolabstract}
\vspace{1em}
]

\section{Introduction}

Visualization is often the first step in interpreting high-dimensional
transcriptomic data. Classical linear methods such as singular value
decomposition and principal component analysis have long been used to summarize
large expression matrices \citep{Alter_2000}. More recent nonlinear methods,
including diffusion maps \citep{Coifman_2005}, UMAP \citep{McInnes_2018}, and
PHATE \citep{Moon_2019}, have become standard tools for displaying cell states,
developmental continua, and branching processes. In single-cell analysis,
visualization is also closely tied to pseudotime inference: methods such as
Monocle \citep{Trapnell_2014}, reversed graph embedding
\citep{Qiu_2017}, and Slingshot \citep{Street_2018} infer ordered or branched
trajectories and then use low-dimensional displays to show progression.

Many biological trajectories, however, are not fundamentally linear. Circadian
time is periodic by definition. The cell cycle is a recurrent process. Some
perturbation and regeneration experiments trace a return toward a starting
state. These settings call for a display in which the primary latent coordinate
is circular. Topological methods address this need by detecting loops and
constructing circular coordinates from persistent cohomology
\citep{de_Silva_2011}. Such methods are attractive because they do not require a
linear ordering and can represent genuinely cyclic sample geometry. They are
also complementary to harmonic analyses of rhythmic transcription, where the
dominant signal is often explicitly modeled by sine and cosine components
\citep{Hughes_2009}.

The method introduced here asks a different but related question. Instead of
first finding a circular coordinate for samples, we ask whether the molecular
features themselves can be ordered on a circle so that the dominant
centered-log-ratio variation looks like a first Fourier mode. This gives a joint
feature-sample visualization. The learned feature circle orders genes or modules
by the phase at which their relative abundance contributes to the cyclic signal.
The sample projection then gives a phase-amplitude display: angle corresponds to
phase, and radius corresponds to strength of the cyclic compositional contrast.

We call the approach \emph{clrcycle}. The compositional step is important because many transcriptomic
measurements are naturally interpreted through relative abundance after library
normalization, and because compositional data are most naturally analyzed using
log-ratio coordinates \citep{Aitchison_1982}. clrcycle imposes a cyclic
feature order and then uses the first circular Fourier mode as the visualization
axis. The result is not intended to replace UMAP, diffusion maps, or topological
loop detection. Rather, it is a targeted projection for cases in which a cyclic
biological coordinate is already plausible and the goal is to obtain an
interpretable sample phase plot together with a feature phase ordering.

This motivation is closely related to simplex-aware visualizations of gene
specificity profiles. When each gene is summarized by weights over a finite set
of biological states, the resulting vectors live on a simplex: raw Euclidean
PCA can distort distances, while barycentric polygon plots can be interpretable
but depend on the ordering of polygon vertices. A useful compromise is to move
through log-ratio coordinates and then align an interpretable polygon to the
observed geometry. The dependence of radial displays on circular coordinate or
anchor order is a known visualization problem in star-glyph and RadViz-style
plots, where the ordering of dimensions around a polygon can strongly affect
the perceived shape and separation of displayed objects \citep{Hu_2021,
Zhu_2022}. clrcycle adopts the same principle for cyclic transcriptomic
structure: use log-ratio geometry first, then learn the polygonal feature order
that makes the two-dimensional display biologically interpretable.

\section{Results}

\subsection{A circadian benchmark in the Hogenesch mouse tissue atlas}

We analyzed the mammalian circadian gene-expression atlas of
\citet{Zhang_2014}. This study sampled 12 mouse tissues every two hours for two
days, giving 24 samples per tissue. Because the sampling times are known, this
dataset provides a direct benchmark: the clrcycle sample angle can be compared with
circadian time, while the sample radius quantifies the strength of the
first-harmonic compositional signal. We used the known times to construct a
supervised reference and then assessed recovery from repeated measurements alone.

\subsection{Supervised reference using time-associated probes}

For a supervised reference, we selected 240 mapped probes with the strongest
first-harmonic association to known circadian time in each tissue, applied the
centered log-ratio transformation, learned a clrcycle feature ordering, and projected
the samples onto the first Fourier mode of the ordered feature circle. The
learned coordinates were aligned to circadian time only for reporting. Because
the same time labels are used for feature selection and evaluation, this analysis
is a descriptive upper reference rather than an independent validation.

In liver, the supervised projection agreed strongly with known circadian time
(\(R=0.951\), mean absolute phase error \(=1.03\) hours; Figure~\ref{fig:liver-supervised}).
Across tissues, agreement ranged from \(R=0.943\) in kidney to \(R=0.983\) in
skeletal muscle, with mean absolute phase errors from 0.59 to 1.16 hours.

The supervised projections exhibit a similar orientation of their elliptical
sample traversals across tissues (Figure~\ref{fig:tissues-supervised}). This
common tilt should not be interpreted as a biological finding: each projection
is independently rotated and, if necessary, reflected to align its angle with
known circadian time. The biologically informative differences are instead the
shape of each traversal and its radius profile. Several peripheral tissues show
a radius maximum near CT8--CT12 and a trough near CT16, whereas skeletal muscle
has its largest radius near CT0 and a trough near CT16. These patterns summarize
tissue-specific changes in the strength of the selected rhythmic compositional
program over the circadian cycle.

The timing of these extrema is biologically plausible, but its interpretation
requires care. In nocturnal mice, CT12 marks the transition into the subjective
night and active phase. Thus, the broad CT8--CT12 radius maxima in several
peripheral tissues are consistent with coordinated transcriptomic remodeling
around that transition. The contrasting CT0 maximum in skeletal muscle is also
credible: tissue-specific circadian phase programs, including markedly divergent
profiles in liver and heart, have been observed in independent transcriptome
studies \citep{Storch_2002,Zhang_2014}. Importantly, a radius maximum is not the
peak time of a particular gene, pathway, or tissue clock. It is a
rotation-invariant summary of the magnitude of the selected 240-probe
compositional contrast, and therefore identifies when that supervised rhythmic
program is most strongly expressed relative to its within-sample background.

\begin{figure*}[t]
\centering
\includegraphics[width=\textwidth]{results/liver_clr_acs_projection.png}
\caption{\textbf{Supervised clrcycle projection of the GSE54650 mouse liver
circadian time course.} Probes were selected by their first-harmonic association
with known circadian time. Left: liver samples projected onto the first circular
Fourier mode after learning a cyclic feature order, colored by circadian time
modulo 24 hours. The dashed circle is the median sample radius. Right: learned
feature circle for the selected probes, with representative clock-associated
genes labeled.}
\label{fig:liver-supervised}
\end{figure*}

\begin{figure*}[t]
\centering
\includegraphics[width=\textwidth]{results/all_tissues/all_tissues_sample_projection_panels.png}
\caption{\textbf{Supervised clrcycle sample projections across the Hogenesch
mouse tissue atlas.} Each panel uses probes selected by first-harmonic
association with known circadian time. Samples are colored by circadian time,
and black lines connect circadian-time centroids. The common direction of the
colored traversal reflects the post-hoc alignment to time; tissue-specific
trajectory shape and radius variation are the meaningful between-tissue
features.}
\label{fig:tissues-supervised}
\end{figure*}

\subsection{Label-free recovery from repeated cycles}

For each tissue, we selected 96 mapped probes whose expression varied
reproducibly between the two sampled circadian cycles, without using circadian
time labels. We then applied the centered log-ratio transformation, learned an
clrcycle feature ordering, and projected the samples onto the first Fourier mode of
the ordered feature circle. Circadian time was used only after fitting to align
and evaluate the display; the arbitrary rotation and reflection of a
two-dimensional circular projection do not otherwise affect the geometry.

In liver, the label-free projection traces the circadian cycle
(Figure~\ref{fig:liver}). Samples progress around the clrcycle plane in circadian
order, and the feature circle arranges selected probes by inferred phase. The
sample angle agrees strongly with known circadian time
(\(R=0.942\), mean absolute phase error \(=1.13\) hours). The radius is not a
confidence interval or fitted orbit; it is the magnitude of the first-harmonic
clr contrast for each sample. Thus large radii indicate times at which the
selected rhythmic log-ratio program is especially strong, while small radii mark
relative troughs in that program.

\begin{figure*}[t]
\centering
\includegraphics[width=\textwidth]{results/all_tissues_unsupervised_repeat_periodic_96/liv_clr_acs_projection.png}
\caption{\textbf{clrcycle projection of the GSE54650 mouse liver circadian time
course after label-free feature selection.} Left: liver samples projected onto
the first circular Fourier mode after learning a cyclic feature order, colored
by circadian time modulo 24 hours for post-hoc validation. The dashed circle is
the median sample radius. Right: learned feature circle for probes selected by
two-cycle repeatability, with representative clock-associated genes labeled.}
\label{fig:liver}
\end{figure*}

The same label-free procedure was run across all 12 tissues
(Figure~\ref{fig:tissues}). Ten tissues had phase agreement \(R>0.89\), ranging
from \(R=0.904\) in aorta to \(R=0.958\) in brainstem, with mean absolute
phase errors from 0.91 to 1.53 hours. Adrenal and brown fat were weaker
(\(R=0.536\) and \(R=0.805\), respectively), illustrating that reproducible
variation alone does not guarantee a dominant circadian program. The radius
profile varied by tissue, suggesting tissue-specific ``tides'' in the strength
of the first-harmonic circadian program.

\begin{figure*}[t]
\centering
\includegraphics[width=\textwidth]{results/all_tissues_unsupervised_repeat_periodic_96/all_tissues_sample_projection_panels.png}
\caption{\textbf{clrcycle sample projections across the Hogenesch mouse tissue
atlas after label-free feature selection.} Each panel shows one tissue from
GSE54650. The 96 probes were selected using reproducibility across the two
sampled cycles, not circadian time. Samples are colored by circadian time only
for post-hoc validation, and black lines connect circadian-time centroids.}
\label{fig:tissues}
\end{figure*}

\begin{table*}[t]
\centering
\caption{\textbf{Summary of circadian recovery across tissues.} Peak and trough
times refer to the largest and smallest clrcycle radii among the two-day samples,
reported modulo 24 circadian hours.}
\label{tab:tissues}
\begin{tabular}{lrrrr}
\toprule
Tissue & \(R\) & Error (h) & Peak CT & Trough CT \\
\midrule
Adrenal & 0.536 & 3.22 & 18 & 14 \\
Aorta & 0.904 & 1.53 & 12 & 16 \\
Brainstem & 0.958 & 0.91 & 12 & 6 \\
Brown fat & 0.805 & 2.04 & 20 & 2 \\
Cerebellum & 0.929 & 1.22 & 12 & 14 \\
Heart & 0.947 & 1.04 & 12 & 18 \\
Hypothalamus & 0.929 & 1.19 & 12 & 2 \\
Kidney & 0.949 & 1.06 & 10 & 16 \\
Liver & 0.942 & 1.13 & 10 & 16 \\
Lung & 0.943 & 1.14 & 10 & 16 \\
Muscle & 0.941 & 1.10 & 0 & 16 \\
White fat & 0.942 & 1.14 & 10 & 16 \\
\bottomrule
\end{tabular}
\end{table*}

\section{Discussion}

clrcycle resolves a cyclic transcriptomic program into two linked objects: a
circular feature ordering and a phase-amplitude projection of samples. In the
circadian atlas, sample angle recovers circadian phase, sample radius measures
the strength of the rhythmic compositional contrast, and feature position
summarizes the phase of the corresponding log-ratio program.

The two analyses separate a reference display from label-free recovery. Selecting
probes by their association with circadian time yields uniformly high agreement,
but uses the evaluation variable during feature selection. Selecting probes by
reproducibility across the two sampled cycles removes this dependence and still
recovers phase in ten tissues. The weaker adrenal and brown-fat projections show
that reproducible variation is not necessarily dominated by a single circadian
mode. This contrast makes the radius and trajectory shape interpretable as
tissue-specific properties of the recovered program rather than consequences of
a uniformly imposed feature set.

The compositional formulation is central to this interpretation. An arbitrary
polygon can make features legible while scrambling their geometry, whereas PCA in
log-ratio coordinates preserves geometry but does not yield phase-readable
features. clrcycle searches for a feature order whose regular-polygon Fourier
coordinates capture clr covariance, thereby joining these two properties. The
ordering is obtained by a discrete heuristic, but its objective is explicit and
its projection is defined in log-ratio coordinates.

clrcycle complements nonlinear embeddings, pseudotime methods, and topological
circular coordinates. Those approaches characterize neighborhoods, branches, or
the existence of a loop in sample space. clrcycle instead tests a specific
representation: whether an ordered set of compositional features supports a
dominant first-harmonic projection. Its contribution is therefore not cyclic
phase inference itself, but a feature-resolved phase-amplitude display for
cyclic transcriptomic structure.

\section{Methods}

\subsection{Input matrix and centered log-ratio transform}

Let \(X\in\mathbb{R}_{\geq 0}^{n\times D}\) be a sample-by-feature matrix after
selecting a feature set. If any entry is nonpositive, we add
\(\alpha=10^{-3}\min\{x_{ij}:x_{ij}>0\}\) to every entry. Each sample is then
transformed by the centered log-ratio map
\[
z_{ij}
= \log x_{ij}
- \frac{1}{D}\sum_{k=1}^{D}\log x_{ik}.
\]
The resulting matrix \(Z\) represents relative log-abundance within the selected
feature set. We center its columns, writing
\(Z_c=Z-\mathbf{1}\bar z^\top\), before computing the feature covariance matrix
\[
\Sigma = \frac{1}{n-1} Z_c^\top Z_c.
\]

\subsection{clrcycle feature ordering}

For a candidate cyclic ordering \(\pi\) of the \(D\) features, place features at
angles
\[
\theta_j = \frac{2\pi j}{D}, \qquad j=0,\ldots,D-1.
\]
Define normalized first-harmonic vectors
\[
c_j = \sqrt{\frac{2}{D}}\cos\theta_j,
\qquad
s_j = \sqrt{\frac{2}{D}}\sin\theta_j.
\]
For ordering \(\pi\), let \(\Sigma_\pi\) be the covariance matrix with rows and
columns permuted by \(\pi\). The clrcycle objective is
\[
F(\pi) = c^\top \Sigma_\pi c + s^\top \Sigma_\pi s.
\]
The learned ordering is
\[
\hat{\pi} = \arg\max_\pi F(\pi).
\]
In the current implementation, this optimization is heuristic. If
\(\lambda_1\geq\lambda_2\) are the two leading eigenvalues and
\(e_1,e_2\) their eigenvectors, we initialize the order by sorting features by
the polar angle of \((\sqrt{\lambda_2}e_{2j},\sqrt{\lambda_1}e_{1j})\). Starting
from this order, we propose 4,000 uniformly sampled pairwise swaps and accept a
swap only when it strictly increases \(F(\pi)\). We use one start per tissue,
with a fixed pseudorandom-number seed.

\subsection{Sample projection}

Given \(\hat{\pi}\), each row \(z^c_i\) of \(Z_c\) is projected onto the first
Fourier mode:
\[
u_i = \sum_{j=1}^{D} z^c_{i,\hat{\pi}_j} c_j,
\qquad
v_i = \sum_{j=1}^{D} z^c_{i,\hat{\pi}_j} s_j.
\]
The polar representation is
\[
\phi_i = \operatorname{atan2}(v_i,u_i),
\qquad
r_i = \sqrt{u_i^2+v_i^2}.
\]
The angle \(\phi_i\) is interpreted as cyclic phase. The radius \(r_i\) is the
strength of the first-harmonic compositional signal for sample \(i\).

\subsection{Summary statistic}

We report
\[
\rho_{\mathrm{clrcycle}}
= \frac{F(\hat{\pi})}{\lambda_1+\lambda_2},
\]
where \(\lambda_1\) and \(\lambda_2\) are the two largest eigenvalues of
\(\Sigma\). This statistic compares the optimized first-harmonic variance with
the maximum variance that any two-dimensional linear projection could capture.
Values near one indicate that the dominant two-dimensional clr covariance is
well represented by a first Fourier mode after cyclic feature ordering.

\subsection{Phase validation}

When known phases are available, write \(\psi_i=\operatorname{atan2}(v_i,u_i)\)
and let \(\phi_i^{\mathrm{ref}}\) be the reference phase. For each of the two
orientations \(\psi_i^{(+)}=\psi_i\) and \(\psi_i^{(-)}=-\psi_i\), we calculate
\[
m_{\pm}=\frac{1}{n}\sum_{i=1}^n
\exp\{\mathrm{i}(\psi_i^{(\pm)}-\phi_i^{\mathrm{ref}})\},
\qquad \delta_{\pm}=\arg(m_{\pm}),
\]
and plot aligned angles \(\hat\phi_i=\psi_i^{(\pm)}-\delta_{\pm}\). We select
the orientation with the largest circular agreement
\[
R=|m_{\pm}|,
\]
breaking an exact tie by smaller mean absolute circular error. Thus the rotation
sets the plotted phase origin but does not change \(R\); only the optional
reflection can change the agreement.
We also report the mean absolute circular phase error, converted to hours for
circadian data. This alignment is used only for validation and plotting; the clrcycle
objective itself is unchanged by global rotation or reflection.

\subsection{Hogenesch atlas analysis}

GSE54650 expression data and GPL6246 probe annotations were downloaded from GEO.
For the supervised reference analysis, we scored positively expressed, mapped
probes by first-harmonic regression against known circadian time. Specifically,
for \(\vartheta_t=2\pi(\mathrm{CT}_t\bmod 24)/24\), we fit
\[
\log x_{tg}-\overline{\log x}_{\cdot g}
= \beta_{cg}\cos\vartheta_t+\beta_{sg}\sin\vartheta_t+\epsilon_{tg}.
\]
Probes were ranked by the fitted fraction of sums of squares
\(R_g^2=\|\hat y_{\cdot g}\|^2/\|y_{\cdot g}\|^2\), with
\(\sqrt{\beta_{cg}^2+\beta_{sg}^2}\) as the secondary criterion, and the top
240 were retained in each tissue. clr transformation, clrcycle ordering, and first
Fourier projection then proceeded as described above. Circadian time was also
used for post-hoc alignment, evaluation, and figure annotation. This analysis
therefore quantifies the best-case display obtained when time-associated probes
are supplied, rather than independent phase recovery.

\subsection{Label-free repeated-cycle analysis}

For each tissue, the 24 samples comprise two cycles with 12 matched sampling
positions per cycle. For every positively expressed, mapped probe, we computed
the log-expression repeatability score
\[
q_g =
\frac{\operatorname{Var}_t\!\left[\bigl(y^{(1)}_{tg}+y^{(2)}_{tg}\bigr)/2\right]}
{\frac{1}{4}\operatorname{Mean}_t\!\left[\bigl(y^{(1)}_{tg}-y^{(2)}_{tg}\bigr)^2\right]+10^{-8}},
\]
where \(y^{(1)}_{tg}\) and \(y^{(2)}_{tg}\) are log expression in the two
cycles at matched sampling position \(t\). The top 96 probes by \(q_g\), with
between-cycle correlation as a secondary ranking criterion, were retained. This
fixed panel size regularizes the feature set uniformly across tissues. Thus
feature selection uses repeated-measurement structure but no circadian-time
values. Each tissue was then analyzed separately by clr transformation, clrcycle
ordering, and first Fourier projection. Circadian time was used only for
post-hoc alignment, evaluation, and figure annotation. The all-tissue figure and
summary table were generated from
\path{results/all_tissues_unsupervised_repeat_periodic_96/}.

\clearpage
\balance
\printbibliography

\end{document}
